% !TeX program = xelatex
% 复制为main.tex，填写config.tex，并按实际参考调整。
\documentclass[12pt,a4paper]{article}
\usepackage{fontspec,xeCJK,geometry,graphicx,amsmath,amssymb}
\usepackage{array,tabularx,booktabs,float,caption,hyperref,titlesec,etoolbox}
\setmainfont{Times New Roman}
\setCJKmainfont[BoldFont=SimHei]{SimSun}
\setCJKsansfont{SimHei}
\setCJKmonofont{SimSun}
\geometry{left=2.5cm,right=2.5cm,top=2.25cm,bottom=2.15cm}
\setlength{\parindent}{2em}
\setlength{\parskip}{0pt}
\renewcommand{\baselinestretch}{1.35}
\renewcommand{\arraystretch}{1.15}
\newcolumntype{Y}{>{\centering\arraybackslash}X}
\newcolumntype{L}{>{\raggedright\arraybackslash}X}
\renewcommand{\figurename}{图}
\renewcommand{\tablename}{表}
\captionsetup{font=small,labelsep=quad,justification=centering}
\titleformat{\section}{\sffamily\fontsize{15pt}{20pt}\selectfont}{\thesection}{.5em}{}
\pagestyle{empty}
\hypersetup{hidelinks}
\providecommand{\ReportTitle}{实验报告}
\providecommand{\ExperimentTitle}{}
\providecommand{\Department}{}
\providecommand{\Grade}{}
\providecommand{\StudentID}{}
\providecommand{\StudentName}{}
\providecommand{\ExperimentDate}{}
\providecommand{\ReportLogo}{}
\InputIfFileExists{config.tex}{}{}
\ifdefined\PRINTVERSION
\newsavebox{\writebox}
\newenvironment{reporttext}{%
\par\begin{lrbox}{\writebox}\begin{minipage}{\textwidth}\setlength{\parindent}{2em}%
\renewcommand{\ref}[1]{0}\renewcommand{\eqref}[1]{(0)}%
}{\end{minipage}\end{lrbox}\vspace*{\dimexpr\ht\writebox+\dp\writebox\relax}\par}
\else
\newenvironment{reporttext}{\par}{\par}
\fi
\raggedbottom
\begin{document}
\begin{titlepage}
\noindent\begin{minipage}[t]{.46\textwidth}
\vspace{0pt}
\ifdefempty{\ReportLogo}{\vspace{2cm}}{%
\IfFileExists{\ReportLogo}{\includegraphics[width=3.2cm]{\ReportLogo}}{\vspace{2cm}}}
\end{minipage}\hfill
\begin{minipage}[t]{.34\textwidth}\vspace{0pt}\raggedleft
\begin{tabular}{|p{4.2cm}|}\hline
\rule{0pt}{1.18cm}\quad 成\quad 绩：\\\hline
\rule{0pt}{1.18cm}\quad 评阅签名：\\\hline
\end{tabular}\end{minipage}
\vspace*{2cm}
\begin{center}{\sffamily\fontsize{23pt}{32pt}\selectfont\ReportTitle}\end{center}
\vspace{1.5cm}
\begin{center}\renewcommand{\arraystretch}{2}
\begin{tabular}{>{\sffamily}r l}
实验题目：&\parbox[c]{9.5cm}{\centering\ExperimentTitle}\\[.3cm]
院系：&\underline{\makebox[9.5cm][c]{\Department}}\\
年级：&\underline{\makebox[9.5cm][c]{\Grade}}\\
学号：&\underline{\makebox[9.5cm][c]{\StudentID}}\\
姓名：&\underline{\makebox[9.5cm][c]{\StudentName}}\\
实验时间：&\underline{\makebox[9.5cm][c]{\ExperimentDate}}\\
\end{tabular}\end{center}
\end{titlepage}
\section{实验目的}
\begin{reporttext}
\end{reporttext}
\section{实验原理}
\begin{reporttext}
\end{reporttext}
\section{实验仪器}
\section{实验内容及步骤}
\begin{reporttext}
\end{reporttext}
\section{数据处理}
\section{实验小结}
\begin{reporttext}
\end{reporttext}
\end{document}
